\chapter[Chapter \thechapter]{}
\lettrine{W}{imsey} presented himself at Mr~Urquhart's house at 9 o'clock the next morning, and found that gentleman at breakfast.

\zz
<I thought I might catch you before you went down to the office,> said his lordship, apologetically. <Thanks awfully, I've had my morning nosebag. No, really, thanks—I never drink before eleven. Bad for the inside.>

<Well, I've found the draft for you,> said Mr~Urquhart pleasantly. <You can cast your eye over it while I drink my coffee, if you'll excuse my going on. It exposes the family skeleton a little, but it's all ancient history now.>

He fetched a sheet of typescript from a side-table and handed it to Wimsey, who noticed, mechanically, that it had been typed on a Woodstock machine, with a chipped lower case p, and an A slightly out of alignment.

<I'd better make quite clear the family connection of the Boyes's and the Urquharts's,> he went on, returning to the breakfast-table, <so that you will understand the will. The common ancestor is old John Hubbard, a highly respectable banker at the beginning of the last century. He lived in Nottingham, and the bank, as usual in those days, was a private, family concern. He had three daughters, Jane, Mary and Rosanna. He educated them well, and they ought to have been heiresses in a mild way, but the old boy made the usual mistakes, speculated unwisely, allowed his clients too much rope—the old story. The bank broke, and the daughters were left penniless. The eldest, Jane, married a man called Henry Brown. He was a schoolmaster and very poor and quite repellantly moral. They had one daughter, Julia, who eventually married a curate, the Rev. Arthur Boyes, and was the mother of Philip Boyes. The second daughter, Mary, did rather better financially, though socially she married beneath her. She accepted the hand of one Josiah Urquhart, who was engaged in the lace-trade. This was a blow to the old people, but Josiah came originally of a fairly decent family, and was a most worthy person, so they made the best of it. Mary had a son, Charles Urquhart, who contrived to break away from the degrading associations of trade. He entered a solicitor's office, did well, and finally became a partner in the firm. He was my father, and I am his successor in the legal business.

The third daughter, Rosanna, was made of different stuff. She was very beautiful, a remarkably fine singer, a graceful dancer and altogether a particularly attractive and spoilt young person. To the horror of her parents she ran away and went on the stage. They erased her name from the family Bible. She determined to justify their worst suspicions. She became the spoilt darling of fashionable London. Under her stage name of Cremorna Garden, she went from one disreputable triumph to another. And, mind you, she had brains—nothing of the Nell Gwyn business about her. She was the take-it-and-keep-it sort. She took everything—money, jewels, apartments meublés, horses, carriages, all the rest of it, and turned it into good consolidated funds. She was never prodigal of anything except her person, which she considered to be a sufficient return for all favours, and I daresay it was. I never saw her till she was an old woman, but before she had the stroke which destroyed her brain and body, she still kept the remains of remarkable beauty. She was a shrewd old woman in her way, and grasping. She had those tight little hands, plump and narrow, that give nothing away—except for cash down. You know the sort.

Well, the long and the short of it was that the eldest sister, Jane—the one who married the schoolmaster—would have nothing to do with the family black sheep. She and her husband wrapped themselves up in their virtue and shuddered when they saw the disgraceful name of Cremorna Garden billed outside the Olympic or the Adelphi. They returned her letters unopened and forbade her the house, and the climax was reached when Henry Brown tried to have her turned out of the Church on the occasion of his wife's funeral.

My grand-parents were less strait-laced. They didn't call on her and didn't invite her, but they occasionally took a box for her performances and they sent her a card for their son's wedding, and were polite in a distant kind of way. In consequence, she kept up a civil acquaintance with my father, and eventually put her business into his hands. He took the view that property was property, however acquired, and said that if a lawyer refused to handle dirty money he would have to show half his clients the door.

The old lady never forgot or forgave anything. The very mention of the Brown-Boyes connection made her foam at the mouth. Hence, when she came to make her will, she put in that paragraph you have before you now. I pointed out to her that Philip Boyes had had nothing to do with the persecution, as, indeed, neither had Arthur Boyes, but the old sore rankled still, and she wouldn't hear a word in his favour. So I drew up the will as she wanted it; if I hadn't, somebody else would have done so, you know.>

Wimsey nodded, and gave his attention to the will, which was dated eight years previously. It appointed Norman Urquhart as sole executor, and, after a few legacies to servants and to theatrical charities, it ran as follows:—

<All the rest of my property whatsoever and wheresoever situated I give to my great-nephew Norman Urquhart of Bedford Row Solicitor for his lifetime and at his death to be equally divided among his legitimate issue but if the said Norman Urquhart should decease without legitimate issue the said property to pass to (here followed the names of the charities previously specified). And I make this disposition of my property in token of gratitude for the consideration shown to me by my said great-nephew Norman Urquhart and his father the late Charles Urquhart throughout their lives and to ensure that no part of my property shall come into the hands of my great-nephew Philip Boyes or his descendants. And to this end and to mark my sense of the inhuman treatment meted out to me by the family of the said Philip Boyes I enjoin upon the said Norman Urquhart as my dying wish that he neither give, lend or convey to the said Philip Boyes any part of the income derived from the said property enjoyed by him the said Norman Urquhart during his lifetime nor employ the same to assist the said Philip Boyes in any manner whatsoever.>

<H'm!> said Wimsey, <that's pretty clear, and pretty vindictive.>

<Yes, it is—but what are you to do with old ladies who won't listen to reason? She looked pretty sharply to see that I had got the wording fierce enough before she would put her name to it.>

<It must have depressed Philip Boyes all right,> said Wimsey. <Thank you—I'm glad I've seen that; it makes the suicide theory a good deal more probable.>

In theory it might do so, but the theory did not square as well as Wimsey could have wished with what he had heard about the character of Philip Boyes. Personally, he was inclined to put more faith in the idea that the final interview with Harriet had been the deciding factor in the suicide. But this, too, was not quite satisfactory. He could not believe that Philip had felt that particular kind of affection for Harriet Vane. Perhaps, though, it was merely that he did not want to think well of the man. His emotions were, he feared, clouding his judgment a little.

He went back home and read the proofs of Harriet's novel. Undoubtedly she could write well, but undoubtedly she knew only too much about the administration of arsenic. Moreover, the book was about two artists who lived in Bloomsbury and led an ideal existence, full of love and laughter and poverty, till somebody unkindly poisoned the young man and left the young woman inconsolable and passionately resolved to avenge him. Wimsey ground his teeth and went down to Holloway Gaol, where he very nearly made a jealous exhibition of himself. Fortunately, his sense of humour came to the rescue when he had cross-examined his client to the verge of exhaustion and tears.

<I'm sorry,> he said; <the fact is, I'm most damnably jealous of this fellow Boyes. I oughtn't to be, but I am.>

<That's just it,> said Harriet, <and you always would be.>

<And if I was, I shouldn't be fit to live with. Is that it?>

<You would be very unhappy. Quite apart from all the other drawbacks.>

<But, look here,> said Wimsey, <if you married me I shouldn't be jealous, because then I should know that you really liked me and all that.>

<You think you wouldn't be. But you would.>

<Should I? Oh, surely not. Why should I? It's just the same as if I married a widow. Are all second husbands jealous?>

<I don't know. But it's not quite the same. You'd never really trust me, and we should be wretched.>

<But damn it all,> said Wimsey, <if you would once say you cared a bit about me it would be all right. I should believe that. It's because you won't say it that I imagine all sorts of things.>

<You would go on imagining things in spite of yourself. You couldn't give me a square deal. No man ever does.>

<Never?>

<Well, hardly ever.>

<That would be rotten,> said Wimsey, seriously. <Of course, if I turned out to be that sort of idiot, things would be pretty hopeless. I know what you mean. I knew a bloke once who got that jealous bug. If his wife wasn't always hanging round his neck, he said it showed he meant nothing to her, and if she did express her affection he called her a hypocrite. It got quite impossible, and she ran away with somebody she didn't care twopence for, and he went about saying that he had been right about her all along. But everybody else said it was his own silly fault. It's all very complicated. The advantage seems to be with the person who gets jealous first. Perhaps you could manage to be jealous of me. I wish you would, because it would prove that you took a bit of interest in me. Shall I give you some details of my hideous past?>

<Please don't.>

<Why not?>

<I don't want to know about all the other people.>

<Don't you, by jove? I think that's rather hopeful. I mean, if you just felt like a mother to me, you would be anxious to be helpful and understanding. I loathe being helped and understood. And, after all, there was nothing in any of them—except Barbara, of course.>

<Who was Barbara?> asked Harriet, quickly.

<Oh, a girl. I owe her quite a lot, really,> replied Wimsey, musingly. <When she married the other fellow, I took up sleuthing as a cure for wounded feelings, and it's really been great fun, take it all in all. Dear me, yes—I was very much bowled over that time. I even took a special course in logic for her sake.>

<Good gracious!>

<For the pleasure of repeating <Barbara celarent darii ferio baralipton.> There was a kind of mysterious romantic lilt about the thing which was somehow expressive of passion. Many a moonlight night have I murmured it to the nightingales which haunt the gardens of St~Johns—though, of course, I was a Balliol man myself, but the buildings are adjacent.>

<If anybody ever marries you, it will be for the pleasure of hearing you talk piffle,> said Harriet, severely.

<A humiliating reason, but better than no reason at all.>

<I used to piffle rather well myself,> said Harriet, with tears in her eyes, <but it's got knocked out of me. You know—I was really meant to be a cheerful person—all this gloom and suspicion isn't the real me. But I've lost my nerve, somehow.>

<No wonder, poor kid. But you'll get over it. Just keep on smiling, and leave it to Uncle Peter.>

When Wimsey got home, he found a note awaiting him.

\begin{mail}{}{Dear Lord~Peter,}

As you saw, I got the job. Miss~Climpson sent six of us, all with different stories and testimonials, of course, and Mr~Pond (the head-clerk) engaged me, subject to Mr~Urquhart's approval.

I've only been here a couple of days, so there isn't very much I can tell you about my employer, personally, except that he has a sweet tooth and keeps secret stores of chocolate cream and Turkish delight in his desk, which he surreptitiously munches while he is dictating. He seems pleasant enough.

But there's just one thing. I fancy it would be interesting to investigate his financial activities. I've done a good bit one way and another with stockbroking, you know, and yesterday in his absence I took a call for him which I wasn't meant to hear. It wouldn't have told the ordinary person anything, but it did me, because I knew something about the man at the other end. Find out if Mr~U. had been doing anything with the Megatherium Trust before their big crash.

Further reports when anything turns up.

\closeletter[Yours sincerely,]{Joan Murchison.}
\end{mail}

<Megatherium Trust?> said Wimsey. <That's a nice thing for a respectable solicitor to get mixed up with. I'll ask Freddy Arbuthnot. He's an ass about everything except stocks and shares, but he does understand them, for some ungodly reason.>

He read the letter again, mechanically noting that it was typed on a Woodstock machine, with a chipped lower case p, and a capital A that was out of alignment.

Suddenly he woke up and read it a third time, noticing by no means mechanically, the chipped p and the irregular capital A.

Then he sat down, wrote a line on a sheet of paper, folded it, addressed it to Miss~Murchison and sent Bunter out to post it.

For the first time, in this annoying case, he felt the vague stirring of the waters as a living idea emerged slowly and darkly from the innermost deeps of his mind.